\documentclass[preprint,review,12pt,authoryear]{elsarticle}

\usepackage[T1]{fontenc}
\usepackage[utf8]{inputenc}
\usepackage{booktabs}
\usepackage{array}
\usepackage{enumitem}
\usepackage{microtype}
\usepackage{url}
\usepackage[hidelinks]{hyperref}

\journal{Journal of Systems and Software}
\biboptions{authoryear}

\begin{document}

\begin{frontmatter}

\title{When Vulnerability Metadata Differ: Routing Trade-Offs across Field-Level NVD--GHSA Strategies}

\author{[Author information to be completed]}
\affiliation{organization={[Affiliation to be completed]},
             addressline={[Address to be completed]},
             city={[City]},
             country={[Country]}}

\begin{abstract}
\textbf{[Result-neutral placeholder.]}
The final abstract will report the frozen corpus, the deterministic RQ1 and RQ2
findings, the action-first/reason-second validation protocol, both
analyst-specific outcomes, and the selected positive or boundary branch. No
analyst-backed result is available in this zero draft.
\end{abstract}

\begin{keyword}
vulnerability metadata \sep vulnerability databases \sep discrepancy analysis
\sep maintenance routing \sep abstention \sep trained-analyst judgment
\end{keyword}

\end{frontmatter}

\section{Introduction}
\label{sec:introduction}

Public vulnerability records support remediation prioritization, affected
software identification, and links from advisories to supporting evidence. The
National Vulnerability Database (NVD) and the GitHub Advisory Database (GHSA)
both publish structured records indexed by Common Vulnerabilities and Exposures
(CVE) identifiers \citep{NVDDevelopers2026,GitHubAdvisoryDatabase2026}. Aligned
records can nevertheless encode severity, affected versions, publication time,
or references differently. A difference alone does not identify its
maintenance consequence. Distinct strings can denote the same value, one record
can contain a strict subset of another, timestamps can mark different
publication events, and the visible evidence can be insufficient to identify a
factual conflict.

Earlier studies establish several parts of this problem. VIEM extracted version
information from public vulnerability reports and compared it with NVD
\citep{Dong2019VIEM}. \citet{Croft2022Severity} studied severity inconsistency
across the Firefox vulnerability-reporting lifecycle and examined the effect of
label-source choice on downstream prediction. \citet{Sun2023Aspect} defined
aspect-level discrepancies across heterogeneous reports, while
\citet{Li2025VuldiffFinder} addressed semantic inconsistency in unstructured
vulnerability information. Other work audits NVD quality, estimates source
quality, automates vulnerability curation, or evaluates affected-version tools
\citep{Anwar2022CleaningNVD,Johnson2018CVSS,Okutan2023Curation,Chen2025AffectedVersions}.
These studies cover discrepancy existence, detection, and metadata quality.
The workflow considered here begins after a structured NVD--GHSA field
difference has already been observed and asks how that difference is routed.

Escalating every unequal value collapses representation, incompleteness,
temporal lag, and factual conflict into one action. Field-specific rules avoid
that collapse but can still encode assumptions that have not been validated.
A type-first strategy assigns actions after a deterministic discrepancy status,
but the resulting status cannot also define the action reference used to
evaluate the strategy. The comparison therefore uses action judgments obtained
independently of strategy output, retains uncertainty, and includes a
field-aware comparator rather than relying on raw inequality.

The study uses a frozen corpus of 8,066 CVE-aligned NVD--GHSA record pairs and
four fields: severity, affected versions, publication date, and references. It
first reports a label-free census of 32,264 field instances. It then compares
three frozen routing strategies: a field-aware simple comparator, a current
type-first candidate, and an abstention-aware type-first candidate. A separate
V3.1 protocol is prepared to evaluate the strategies with maintenance actions
assigned by two independent doctoral-student trained analysts. The protocol
collects actions before discrepancy reasons so that the discrepancy taxonomy
does not directly define the action reference.

The research questions are:

\begin{description}[leftmargin=1.7cm,style=nextline]
\item[RQ1] Across 8,066 CVE-aligned NVD--GHSA record pairs, how do deterministic
field statuses distribute for severity, affected versions, publication date,
and references?
\item[RQ2] How do a field-aware simple comparator, a current type-first
candidate, and an abstention-aware type-first candidate allocate field
instances across actions, and where do their conflict-escalation, abstention,
and total manual-route outputs differ?
\item[RQ3] When two independent trained analysts assign maintenance actions to
the same frozen formal cases, do their judgments differentiate the three
routing strategies in a consistent direction, and where do reliability,
agreement, coverage, abstention, or shared-miss boundaries emerge?
\end{description}

The result-neutral thesis is that three frozen routing strategies produce
different conflict-escalation, abstention, and total manual-route allocations
on four fields of CVE-aligned NVD--GHSA records. Independent trained-analyst
judgments are reserved for testing whether those deterministic differences
correspond to differentiated maintenance actions or expose an empirical
decision boundary.

The study makes three bounded contributions:

\begin{enumerate}
\item a reproducible, snapshot-bounded census of 8,066 CVE-aligned record pairs
and 32,264 field instances, with deterministic statuses for four fields;
\item a three-strategy routing comparison with separate conflict-escalation,
abstention, and total-manual-route accounting; and
\item a frozen two-analyst evaluation that can support either consistent
strategy differentiation by both analysts or a reported reliability,
agreement, coverage, abstention, shared-miss, or identifiability boundary.
\end{enumerate}

RQ1 and RQ2 are supported by deterministic evidence. RQ3 retains a positive and
a boundary branch pending valid analyst returns. The positive branch requires
both analysts to clear the frozen reliability, paired-direction, event-floor,
coverage, and manual-loss gates. Otherwise, the boundary branch reports the
failed gates, disagreements, abstentions, and uncertain outcomes.

\section{Background and Task Definition}
\label{sec:background}

\subsection{Research object}

The unit of alignment is a CVE identifier shared by an NVD record and a reviewed
GHSA record under the frozen data pipeline. The unit of analysis is a field
instance: a CVE, a field, the NVD value, and the GHSA value. CVE alignment is an
upstream identifier relation. It is not independent semantic adjudication and
does not cover records without a usable CVE identifier, ambiguous multi-CVE
mappings, or true no-match cases.

The four primary fields are severity, affected versions, publication date, and
references. CWE identifiers are outside the low-human V3.1 study. Each primary
field has a different identity contract. Severity comparison distinguishes
missing scores, representation differences, and materially different values.
Affected-version comparison depends on package or product identity and range
semantics. Publication dates can refer to different source events. Reference
comparison is set- and URL-sensitive and can expose aliases, source additions,
or inaccessible evidence.

\subsection{Observation, reason, and action}

The deterministic pipeline assigns one of five statuses when its field-specific
rules permit: equivalent (EQ), representation discrepancy (RD), incomplete
(INC), temporal discrepancy (TD), or factual conflict (FC). The analyst protocol
also permits \emph{uncertain}. These labels describe candidate reasons for an
observed difference; they are not source-truth labels.

Maintenance action is recorded separately as no action, enrich record, wait for
sync, conflict escalation, or abstain. Conflict escalation places a case in a
conflict queue. Abstention also routes a case to manual review without asserting
that a factual conflict exists. Total manual route is therefore defined as
conflict escalation plus abstention. This quantity counts strategy outputs; it
does not measure elapsed labor, monetary cost, or observed operational workload.

\subsection{Separation from source adjudication}

The task asks what should happen next under the visible record context. It does
not require an analyst to declare either source globally authoritative.
Abstention and uncertainty remain admissible when the visible material cannot
support an action or reason. Earlier evidence-driven source-adjudication
experiments are retained only as bounded failure evidence because their
non-human cohorts and abstentions cannot supply the independent action reference
required for RQ3.

\section{Related Work}
\label{sec:related}

\subsection{Vulnerability-report and database discrepancies}

\citet{Dong2019VIEM} extracted software-version information from public reports
and compared the extracted values with NVD. \citet{Croft2022Severity} measured
severity disagreement across Bugzilla, Mozilla advisories, and NVD, then
examined how label-source choice affected prediction. These studies motivate
field-specific analysis, but neither supplies four-field NVD--GHSA maintenance
actions.

\citet{Sun2023Aspect} provide the closest task-level overlap: their study aligns
heterogeneous reports, extracts seven vulnerability aspects, and analyzes
aspect-level discrepancy types and detection. VuldiffFinder also treats
unstructured semantic inconsistency as a detection task
\citep{Li2025VuldiffFinder}. The present task concerns the next workflow step:
routing differences between structured NVD and GHSA records to maintenance
actions. Actions precede reasons in the analyst protocol, the main comparator is
field-aware, and abstention is included in total manual routing.

\subsection{Metadata quality, curation, and affected versions}

NVD-focused work has assessed completeness and consistency, proposed automatic
corrections, and modeled uncertainty in severity sources
\citep{Anwar2022CleaningNVD,Johnson2018CVSS}. Automated curation maps text to
structured attributes and evaluates the resulting predictions
\citep{Okutan2023Curation}. These studies evaluate metadata or attribute
generation rather than the routing of a difference between two structured
sources.

Affected-version analysis requires release histories, branches, backports, and
artifact identity. \citet{Chen2025AffectedVersions} benchmarked 12 tools on
1,128 C/C++ vulnerabilities under a common task contract. It is the closest
same-field benchmark identified in the literature audit. Its task is
affected-version identification rather than cross-source routing, and the tools
have not been reproduced on the present corpus. The study therefore defines an
external task boundary; it is not reported as a performance baseline for the
present field rules.

\subsection{Selective prediction and human routing}

Selective prediction and learning to defer formalize the choice between an
automatic prediction and routing to another decision maker.
\citet{Mozannar2020Defer} jointly learn a classifier and rejector from labels and
expert decisions under a system loss. The analogy is useful because deferral
depends on the automatic component and the downstream decision maker. The
present strategies are deterministic and frozen before analyst exposure; no
rejector is learned. Learning to defer is therefore a theoretical neighbor, not
a reproduced same-task baseline.

\subsection{Public resources and GHSA process}

VulZoo aggregates multiple vulnerability-intelligence sources
\citep{Ruan2024VulZoo}; CVEfixes links CVEs to fixes and code
\citep{Bhandari2021CVEfixes}; VFCFinder pairs advisories with candidate patches
\citep{Dunlap2024VFCFinder}; and data-quality studies audit vulnerability
datasets and their construction \citep{Croft2023DataQuality}. These resources
can support aggregation, repair linkage, and quality auditing. They do not
provide same-contract NVD--GHSA maintenance-action labels or an independent
correctness oracle. The broader vulnerability-database literature also treats
data sources, schemas, quality, and downstream uses as distinct concerns
\citep{Li2023Anatomy}.

\citet{Segal2026GHSA} characterize the GHSA review pipeline, identify two
review-latency regimes, and release data and analysis code. That evidence
provides process context for publication-date differences. It does not establish
field correctness or temporal generalization for the frozen snapshot used here.

\subsection{Positioning}

Discrepancy detection, metadata-quality assessment, attribute generation,
affected-version benchmarking, and human deferral are established but separate
tasks. The present contract combines four structured fields, a maintenance
action vocabulary, three frozen routing strategies, and two independent
action-first/reason-second trained-analyst passes. Its empirical differential is
action-oriented routing with explicit abstention and a field-aware comparator;
it is not a claim to be the first discrepancy taxonomy.

\section{Corpus and Deterministic Analysis}
\label{sec:corpus}

\subsection{Data sources and alignment}

The frozen pipeline normalizes NVD and reviewed GHSA records and aligns them by
CVE identifier. The resulting field view contains 8,066 aligned record pairs
and four fields per pair, yielding 32,264 field instances. The field-view
SHA-256 is:
\begin{center}
\texttt{c4bb405399bd0050c206b63ece95771c4}\\
\texttt{f566a9a3968f123a16cf02a3b5cc3a2}.
\end{center}
The corpus supports statements about this snapshot and pipeline. It is not a
probability sample of all vulnerability records, and CVE alignment excludes
open-world matching cases.

\subsection{Frozen strategies and reference arms}

The field-aware simple comparator maps field-specific conditions directly to
maintenance actions. The current type-first candidate maps the deterministic
status to an action and does not explicitly abstain. The abstention-aware
type-first candidate routes selected conditions with insufficient
identifiability to abstention.

Four additional arms bound the design. Raw and canonical non-equality provide
binary lower references. An always-manual arm sends every instance to conflict
escalation, and an abstain-all arm sends every instance to abstention. These
four arms are reference cases rather than main same-contract competitors.

\subsection{Reproducibility and sampling use}

The label-free analyzer and an independent verifier bind the field view,
strategy implementation, counts, and sampling capacity. Their agreement
establishes mechanical reproducibility of the deterministic outputs; it does not
establish the semantic validity of the rules. The census was also used to check
whether a fixed analyst sample could contain strategy-disagreement and
shared-no-manual cases without changing the study objective.

\section{Analyst-Bounded Validation Protocol}
\label{sec:protocol}

\subsection{Analyst role and bounded calibration}

The protocol specifies two independent doctoral students trained as analysts
for this study, denoted Reviewer A and Reviewer B in the artifacts. It requires
independent work without AI assistance. The construct is trained-analyst
judgment; practitioner behavior is outside its scope.

V3.1 contains 20 calibration-1 cases, a CVE-disjoint 20-case calibration-2
reserve, and 120 formal cases. Calibration-2 is opened only if calibration-1
action agreement is below 0.60 or the guideline changes materially. A second
action agreement below 0.60 terminates formal distribution. The formal
allocation is 50 severity, 50 affected-version, 10 publication-date, and 10
reference cases. Cases are unique within phases, and the phase CVE sets are
pairwise disjoint.

\subsection{Action first, reason second}

Each analyst first assigns an action using the visible record pair and allowed
context. The return is validated and the action stage is hash-locked before a
reason packet is released. The analyst then assigns EQ, RD, INC, TD, FC, or
uncertain. This order prevents the project's reason taxonomy from directly
defining the action used for strategy evaluation. The same person sees the same
case twice, so same-analyst action--reason association is treated as an upper
bound rather than the primary association.

\subsection{Blinding and file governance}

Analyst-visible objects use recursive allowlists. Strategy outputs,
deterministic statuses, selection cells, weights, AI candidates, prior reviews,
the other analyst's materials, and future-stage packets are excluded. URLs can
reveal source identity, so the protocol is strategy-blinded but not perfectly
source-blinded.

\subsection{Estimands and stop rules}

Pre-adjudication analysis reports analyst-specific raw action agreement,
nominal Krippendorff alpha, uncertainty, disagreement matrices, and reason
agreement as RQ3 validity diagnostics. Cross-analyst action--reason association
is primary; same-analyst association is secondary. Paired action matches on
strategy-disagreement rows, exact McNemar discordance, blocked intervals,
manual-route coverage, abstention, and design-weighted sensitivity are computed
separately for each analyst.

The formal set contains a fixed 34-case shared-no-manual audit: 15 severity and
19 affected-version cases. An analyst conflict-escalation action on one of these
cases identifies a shared miss within the sample. The inference remains limited
to those 34 cases. The positive branch requires each analyst to provide at least
29 conflict actions, no lower type-first manual coverage, a one-sided
simple-only-loss upper bound below \(\delta_{\mathrm{manual}}=0.10\), the frozen
paired direction, and no systematic contradictory failure. A failure for either
analyst selects the boundary branch.

\subsection{Adjudication}

Author adjudication, if performed, is strategy-blinded, secondary, and reported
after the pre-adjudication results. Sensitivity analysis excludes every
adjudicated case. Adjudication cannot repair low independent agreement or
reverse a failed strategy gate.

\section{Results}
\label{sec:results}

\subsection{RQ1: Deterministic discrepancy landscape}

% Generated by experiments/paper_artifacts/build_jss_deterministic_tables.py
% Source: results/jss/t1_routing_precheck_v1/analysis.json
\begin{table}[t]
\centering
\caption{Deterministic discrepancy-status counts by field.}
\label{tab:rq1-status}
\footnotesize
\setlength{\tabcolsep}{3pt}
\begin{tabular}{lrrrrrr}
\toprule
Field & EQ & RD & INC & TD & FC & Total \\
\midrule
Severity & 3106 & 3178 & 33 & 0 & 1749 & 8066 \\
Affected versions & 425 & 3936 & 3054 & 0 & 651 & 8066 \\
Publication date & 0 & 6169 & 0 & 1897 & 0 & 8066 \\
References & 0 & 300 & 7763 & 0 & 3 & 8066 \\
\bottomrule
\end{tabular}
\begin{minipage}{0.98\linewidth}
\footnotesize\textit{Note:} EQ = equivalent; RD = representation discrepancy; INC = incomplete; TD = temporal discrepancy; FC = factual conflict. Counts are deterministic rule outputs, not human-verified factual labels.
\end{minipage}
\end{table}


Table~\ref{tab:rq1-status} reports 8,066 observations for each field. Severity
contains 3,106 equivalent, 3,178 representation-discrepancy, 33 incomplete, and
1,749 factual-conflict outputs. Affected versions contain 425 equivalent, 3,936
representation-discrepancy, 3,054 incomplete, and 651 factual-conflict outputs.
Publication date contains 6,169 representation and 1,897 temporal discrepancies.
References contain 300 representation discrepancies, 7,763 incomplete
instances, and three factual-conflict outputs. These counts are deterministic
rule outputs, not verified factual labels or database-quality measurements.

\textbf{Answer to RQ1.} Status allocations differ by field in the frozen census.
Representation discrepancies have the largest count for severity and publication
date, incomplete values have the largest count for references, and
affected-version instances are distributed across equivalent, representation,
incomplete, and factual-conflict statuses.

\subsection{RQ2: Deterministic routing comparison}

% Generated by experiments/paper_artifacts/build_jss_deterministic_tables.py
% Source: results/jss/t1_routing_precheck_v1/analysis.json
\begin{table}[t]
\centering
\caption{Deterministic routing allocations by strategy.}
\label{tab:rq2-actions}
\footnotesize
\setlength{\tabcolsep}{3pt}
\begin{tabular}{lrrrrrrr}
\toprule
Strategy & No act. & Enrich & Wait & Conflict & Abstain & Manual & Total \\
\midrule
Field-aware & 15182 & 12009 & 1897 & 1780 & 1396 & 3176 & 32264 \\
Type-first current & 17114 & 10850 & 1897 & 2403 & 0 & 2403 & 32264 \\
Type-first abstention & 15465 & 10776 & 1897 & 1706 & 2420 & 4126 & 32264 \\
\bottomrule
\end{tabular}
\begin{minipage}{0.98\linewidth}
\footnotesize\textit{Note:} Manual total is Conflict + Abstain. Counts describe strategy outputs on 32,264 field instances; they do not measure labor, correctness, safety, or utility.
\end{minipage}
\end{table}

% Generated by experiments/paper_artifacts/build_jss_deterministic_tables.py
% Source: results/jss/t1_routing_precheck_v1/analysis.json
\begin{table}[t]
\centering
\caption{Pairwise action disagreements across routing strategies.}
\label{tab:rq2-disagreements}
\footnotesize
\setlength{\tabcolsep}{3pt}
\begin{tabular}{lrrrrr}
\toprule
Pair & Sev. & Affected & Date & Refs & Total \\
\midrule
Simple/current & 5 & 2247 & 0 & 303 & 2555 \\
Simple/abstention & 263 & 1766 & 0 & 303 & 2332 \\
Current/abstention & 258 & 2159 & 0 & 3 & 2420 \\
\bottomrule
\end{tabular}
\begin{minipage}{0.98\linewidth}
\footnotesize\textit{Note:} Each cell counts field instances on which the two named frozen strategies emit different actions. Zero is an observed census count, not missing data.
\end{minipage}
\end{table}


Table~\ref{tab:rq2-actions} reports action allocations for all three main
strategies. The field-aware simple comparator routes 3,176 instances manually,
the current type-first candidate routes 2,403, and the abstention-aware
type-first candidate routes 4,126. These are strategy-output counts; the current
type-first candidate has no abstention action under the frozen mapping.

The three pairwise comparisons in Table~\ref{tab:rq2-disagreements} differ on
2,555, 2,332, and 2,420 field instances, respectively. For the focal
field-aware-simple versus abstention-aware comparison, the 2,332 action
differences comprise 263 severity, 1,766 affected-version, no publication-date,
and 303 reference instances. The abstention-aware strategy emits 1,706 conflict
escalations compared with 1,780 for the field-aware comparator, a difference of
\(-74\). Once abstentions are included, its total manual route is 4,126 compared
with 3,176, a difference of \(+950\).

\textbf{Answer to RQ2.} The frozen strategies allocate the same 32,264 field
instances differently. Most disagreement in the focal comparison occurs in
affected versions, and the two strategies agree on publication-date actions.
A smaller conflict queue coexists with a larger total manual route when
abstention is counted. The census does not measure labor, safety, correctness,
or operational utility.

\subsection{RQ3: Analyst-bounded validation}

\textbf{[REAL-HUMAN RESULTS PLACEHOLDER.]}
No RQ3 result is available. This subsection must not be populated with
AI-generated, synthetic, calibration-only, or author-imputed labels.

After valid returns, this subsection will report both analysts separately:
action distributions including abstention; raw agreement and nominal
Krippendorff alpha overall and by field; reason distributions, uncertainty, and
disagreement matrices; cross-analyst and same-analyst action--reason
associations; the calibration path used; paired action-match differences and
exact discordance; conflict, abstention, and total-manual-route coverage; the
34-case shared-no-manual audit; the reviewer-specific event and
\(\delta_{\mathrm{manual}}\) gates; weighted sensitivity; and an
adjudication-exclusion sensitivity.

The positive branch is admissible only if both analysts clear every frozen
gate and support the same direction. Otherwise, the boundary branch reports the
observed reliability, paired-direction, event-floor, coverage, manual-loss, or
identifiability boundary without changing fields, cases, strategies,
thresholds, or labels.

\section{Discussion}
\label{sec:discussion}

\subsection{What the deterministic results establish}

RQ1 describes the snapshot-bounded allocation of rule-based discrepancy
statuses. RQ2 shows that conflict escalation and total manual routing can move
in opposite directions once abstention is included. Reporting both quantities
keeps abstentions visible in the routing comparison. These results establish
allocation differences; they do not identify which strategy is more accurate
or useful.

\subsection{Positive-branch interpretation}

\textbf{[CONDITIONAL TEXT.]}
If both analysts clear the frozen gates, the final discussion will interpret the
result as a bounded comparison under independent trained-analyst actions. It
will report which fields carry the observed differentiation and whether
abstention changes coverage under the frozen contract. The interpretation will
remain about action alignment and manual-route coverage, not source correctness,
maintenance effort, or deployment safety.

\subsection{Boundary-branch interpretation}

\textbf{[CONDITIONAL TEXT.]}
If a gate fails, the final discussion will retain the disagreement or
insufficient event count as an empirical result. Any interpretation---for
example, insufficient visible information, field-dependent ambiguity, unstable
action semantics, or a shared blind spot---will be selected only when the
observed matrices and cases support it. Failure of the positive branch will not
be generalized to all vulnerability-metadata routing.

\subsection{Limits of automated reconciliation}

Historical non-human adjudication cohorts in the repository were frequently
unresolved or dependent on available evidence under the tested protocols.
Those results explain the retained abstention action. They are not merged with
V3.1 analyst outcomes, presented as a successful adjudication method, or used to
rank NVD and GHSA.

\section{Threats to Validity}
\label{sec:threats}

\paragraph{Construct validity.}
Maintenance actions are study constructs rather than logs of production curator
behavior. The analysts are trained for this study and are not represented as
practitioners. Action-first ordering reduces direct taxonomy circularity but
does not eliminate repeated-case anchoring. URLs may reveal source identity.
Uncertainty and pre-adjudication agreement are reported, and low calibration or
formal reliability blocks a positive strategy claim.

\paragraph{Internal validity.}
The deterministic field rules can be semantically wrong even when their
implementation is reproducible. Sampling is stratified around strategy
comparison and contains a fixed shared-no-manual audit; it is not an
unqualified random sample of all instances. Strategies, sampling cells,
thresholds, and analyses were frozen before analyst exposure. Analyst-specific
paired comparisons reduce dependence on pooled judgments but remain conditional
on the action and loss contract.

\paragraph{Conclusion validity.}
A limited number of analyst conflict-escalation actions may leave the positive
branch statistically unidentified. The frozen 25/29 reporting thresholds and
one-sided manual-loss upper bound control the permitted interpretation; they do
not guarantee a positive result or operational safety. Weighted sensitivity can
have low effective sample size and does not replace the primary paired
comparison. Field-level summaries remain descriptive unless a confirmatory
hierarchy or correction was frozen.

\paragraph{External validity.}
The corpus contains CVE-aligned NVD and reviewed GHSA records from one snapshot
and pipeline. It excludes unmatched records, records without a CVE, ambiguous
multi-CVE mappings, and broader advisory sources. Results may not transfer to
future snapshots, unreviewed GHSAs, other data sources, or practitioner
workflows. Existing snapshot-external data do not form a bilateral post-freeze
temporal cohort.

\paragraph{Reproducibility and artifact validity.}
Hashes, manifests, validators, and independent mechanical verifiers bind files
and computations. Their scope is mechanical reproducibility. Analyst returns
can contain private or sensitive information; any public artifact must separate
anonymized analysis data from analyst-private material and follow the final
author and institutional disposition.

\section{Conclusion}
\label{sec:conclusion}

\textbf{[CONCLUSION PLACEHOLDER.]}
The final conclusion will restate the bounded RQ1 and RQ2 findings and report
either the analyst-consistent positive result or the preserved RQ3 boundary.

\section*{Data availability}

\textbf{[AUTHOR TO COMPLETE.]}
The final statement must provide the public repository and persistent
identifier for shareable research data, or explain why particular data cannot
be shared. Private analyst materials must not be included in a public deposit.

\section*{Code availability}

\textbf{[AUTHOR TO COMPLETE.]}
The final statement must identify an archival software version, persistent
identifier, and license.

\section*{Funding}

\textbf{[AUTHOR TO COMPLETE.]}

\section*{Declaration of competing interest}

\textbf{[AUTHOR TO COMPLETE.]}

\section*{CRediT authorship contribution statement}

\textbf{[AUTHOR TO COMPLETE.]}

\section*{Human participation and ethics}

\textbf{[AUTHOR/INSTITUTION TO COMPLETE.]}
The final statement must report the applicable institutional determination,
recruitment, consent, compensation, and conflict conditions accurately.

\section*{Declaration of generative AI and AI-assisted technologies in the writing process}

\textbf{[AUTHOR TO REVIEW AND COMPLETE USING THE CURRENT JOURNAL WORDING.]}
Codex assisted preparation and editing of this zero draft. The final declaration
must identify the tool and purpose, describe author review and editing, and
state that the authors retain responsibility for the article.

\bibliographystyle{elsarticle-harv}
\bibliography{references}

\end{document}
